\documentclass[11pt, a4paper]{article}
\usepackage[left=25mm, right=20mm, top=20mm, bottom=20mm]{geometry}
\setlength{\parindent}{0pt}
\setlength{\parskip}{0.5em} 

\begin{document}

\begin{titlepage}
    \centering
    \vspace*{1cm}
    {\large University of Primorska\par}
    \vspace{0.2cm}
    {\large Faculty of Mathematics, Natural Sciences and Information Technologies\par}

    \vspace{4cm}

    {\Huge \bfseries Overboard Code Cracker\par}

    \vspace{4.5cm}

    \begin{flushright}
        {\large Authors: Sergei Savinkov, Artem Nesterov, Ivan Mikhailov\par
        Course: Programsko inženirstvo EN\par
        Group: 01\par}
    \end{flushright}

    \vfill

    {\large Koper, 2026\par}

\end{titlepage}

%Homework 1
\section{Problem definition}
\subsection{Problem}
Two elderly friends now live in different cities but they have a long tradition of playing logic-based Code Cracker games remotely.
They set up video cameras over their game boards to track each others attempts.
However, due to the cognitive load of managing intermediate guessing rows, 
tracking historical turn sequences and manually computing exact feedback for each row, 
they frequently lose count, misplace pegs or make calculation errors.
This leads to confusion, mutual accusations of cheating and ultimately undermines their desire to continue playing together.

\subsection{Proposed solution}
A software system that connects two physical Code Cracker boards located in different cities. Each player uses a camera positioned above their board. 
Computer vision identifies the placement of coloured pegs and detects when a new guess has been completed. 
When a player places new pegs the system detects this, converts everything into a digital representation and transmits it to the other participant in order to set the new game state. 
The receiving application then compares the expected game state with the board state detected by its own camera.

This allows both physical boards to be tracked in the same digital game history. 
If a board does not reflect the expected state or if the feedback pegs violate the rules of code cracker the system alerts both players and allows them to review the inconsistency. 
Players continue using their physical boards, while the software coordinates the game state between their.

\subsection{Requirements the solution must meet}
\begin{itemize}
    \item The system must support two players playing Code Cracker remotely using separate physical boards and cameras.
    \item The system must use computer vision to recognize coloured pegs and completed guesses on each board.
    \item The system must update a shared digital game state with recognized guesses, which maintains the shared game history and synchronizes updates between both players.
    \item The system must compare each physical board with its expected digital state and detect misplaced, missing or unexpected pegs.
    \item The system must recognize the two types of feedback pegs and verify that they are consistent with the rules of Code Cracker.
    \item The system must notify both players of inconsistencies and allow them to review previous guesses and correct recognition errors.
\end{itemize}

\end{document}